\chapter{Reduction to finitely many sizes}

\begin{definition}\leanok
[Nonstandard counterexample at a size]
  \label{def:counterexample_at}
  \lean{Ninegates.CounterexampleAt}
  \uses{def:live_wait,def:standard}
  % SOURCE: half-formal.qmd, “把任意反例压缩到有限规模” (read from references/half-formal.qmd)
  % SOURCE QUOTE: “称一个非空听的 $n$ 面听牌型为**非标准牌型**，如果它不等于 $G_n=(3,1,\ldots,1,3)$。”
  \textit{Source: half-formal.qmd, finite-size reduction section.}
  A counterexample at size \(n\) is a count vector \(c\) that is a live
  all-rank wait and is not equal to \(G_n\).
\end{definition}

\begin{definition}\leanok
[The finite set of small sizes]
  \label{def:small_sizes}
  \lean{Ninegates.smallSizes}
  % SOURCE: half-formal.qmd, #lem-bound (read from references/half-formal.qmd)
  % SOURCE QUOTE: “那么存在 $n_0\in\{6,9,12,15,18,21\}$，以及一个非标准的非空听 $n_0$ 面听牌型。”
  \textit{Source: half-formal.qmd, Lemma \texttt{\#lem-bound}.}
  Define
  \[
    \mathcal S=\{6,9,12,15,18,21\}.
  \]
\end{definition}

\begin{definition}

\leanok
[Internal-coordinate embedding]
  \label{def:internal_embedding}
  \lean{Ninegates.internalEmbedding}
  For \(n\) ranks, the internal coordinate indexed by \(r\in\{0,\ldots,n-3\}\)
  is the original coordinate \(r+1\).  This gives an injective map from
  \(n-2\) auxiliary indices into the coordinates of the original hand.
\end{definition}

\begin{definition}

\leanok
[The finite set of internal coordinates]
  \label{def:internal_indices}
  \lean{Ninegates.internalIndices}
  \uses{def:internal_embedding}
  The set of internal coordinates is the image of the preceding embedding.
  Equivalently, it is the finite set of coordinates strictly between the two
  endpoints.
\end{definition}

\begin{lemma}

\leanok
[Membership in the internal-coordinate set]
  \label{lem:mem_internal_indices}
  \lean{Ninegates.mem_internalIndices}
  \uses{def:internal_indices,def:internal_embedding}
  A coordinate \(i\) belongs to the internal-coordinate set exactly when
  \(1\leq i\) and \(i+1<n\), in zero-based indexing.
\end{lemma}
\begin{proof}

\leanok
  If \(i\) is in the image, then \(i=r+1\), so it is not the first coordinate
  and the bound on \(r\) shows it is not the last.  Conversely, from
  \(1\leq i\) and \(i+1<n\), the auxiliary index \(i-1\) lies in
  \(\{0,\ldots,n-3\}\) and maps back to \(i\).
\end{proof}

\begin{lemma}

\leanok
[Cardinality of the internal-coordinate set]
  \label{lem:card_internal_indices}
  \lean{Ninegates.card_internalIndices}
  \uses{def:internal_indices,def:internal_embedding}
  The number of internal coordinates is \(n-2\).
\end{lemma}
\begin{proof}

\leanok
  The internal-coordinate set is the image of an injective map from a finite
  type with \(n-2\) elements, so its cardinality is \(n-2\).
\end{proof}

\begin{lemma}

\leanok
[Small admissible sizes are exactly the listed sizes below twenty-three]
  \label{lem:small_sizes_mem_of_admissible_le_22}
  \lean{Ninegates.smallSizes_mem_of_admissible_le_22}
  \uses{def:small_sizes}
  If \(6\leq n\), \(3\mid n\), and \(n\leq22\), then
  \(n\in\mathcal S=\{6,9,12,15,18,21\}\).
\end{lemma}
\begin{proof}

\leanok
  Write \(n=3k\).  The inequalities give \(2\leq k\leq7\), and checking these
  six possible values of \(k\) gives precisely the six elements of
  \(\mathcal S\).
\end{proof}

\begin{lemma}

\leanok
[The total internal excess is at most six]
  \label{lem:internal_excess_le_six}
  \lean{Ninegates.internal_excess_le_six}
  \uses{def:internal_indices,lem:mem_internal_indices,lem:card_internal_indices,def:live_wait,def:total,thm:interior_positive}
  Let \(c\) be a live all-rank wait at an admissible size \(n\).  Then
  \[
    \sum_{\text{internal }i}(c_i-1)\leq6.
  \]
\end{lemma}
\begin{proof}

\leanok
  Interior positivity gives \(1\leq c_i\) on every internal coordinate, hence
  \((c_i-1)+1=c_i\) there.  Therefore
  \[
    \sum_{\text{internal }i}(c_i-1)+(n-2)
      =\sum_{\text{internal }i}c_i.
  \]
  The internal sum is bounded by the total tile count, which is \(n+4\) for a
  live all-rank wait.  Subtracting the \(n-2\) internal coordinates gives the
  bound \(6\).
\end{proof}

\begin{lemma}

\leanok
[Exceptional internal coordinates are bounded by internal excess]
  \label{lem:exception_card_le_internal_excess}
  \lean{Ninegates.exception_card_le_internal_excess}
  \uses{def:internal_indices}
  Suppose every internal coordinate of \(c\) has multiplicity at least \(1\).
  Then the number of internal coordinates with \(c_i\ne1\) is at most
  \[
    \sum_{\text{internal }i}(c_i-1).
  \]
\end{lemma}
\begin{proof}

\leanok
  Let \(E\) be the internal coordinates where \(c_i\ne1\).  Each \(i\in E\)
  has \(c_i\geq1\) and is not equal to \(1\), so \(c_i-1\geq1\).  Summing over
  \(E\) gives \(|E|\leq\sum_{i\in E}(c_i-1)\), and enlarging the finite set to
  all internal coordinates can only increase the nonnegative sum.
\end{proof}

\begin{lemma}

\leanok
[Seven fixed blocks force seven exceptional coordinates]
  \label{lem:seven_le_exception_card_of_no_three}
  \lean{Ninegates.seven_le_exception_card_of_no_three}
  \uses{def:internal_indices,lem:mem_internal_indices}
  If \(24\leq n\) and there is no internal block of three consecutive
  coordinates all equal to \(1\), then at least seven internal coordinates
  satisfy \(c_i\ne1\).
\end{lemma}
\begin{proof}

\leanok
  Consider the seven disjoint internal blocks
  \([1,4),[4,7),\ldots,[19,22)\) in zero-based coordinates.  Each block is
  contained in the internal coordinates when \(24\leq n\).  Since no block of
  length three is entirely made of \(1\)'s, choose one coordinate in each block
  where \(c_i\ne1\).  The blocks are disjoint, so these seven choices are
  distinct exceptional internal coordinates.
\end{proof}

\begin{theorem}\leanok
[Minimal counterexample size bound]
  \label{thm:counterexample_size_bound}
  \lean{Ninegates.counterexample_size_bound}
  \uses{def:counterexample_at,def:small_sizes,lem:small_sizes_mem_of_admissible_le_22,lem:mem_internal_indices,lem:internal_excess_le_six,lem:exception_card_le_internal_excess,lem:seven_le_exception_card_of_no_three,thm:interior_positive,thm:delete_three_live_wait,thm:delete_three_standard}
  % SOURCE: half-formal.qmd, #lem-bound (read from references/half-formal.qmd)
  % SOURCE QUOTE: “假设存在某个 $n\ge6$，其中 $3\mid n$，并存在一个非标准的非空听 $n$ 面听牌型。那么存在 $n_0\in\{6,9,12,15,18,21\}$，以及一个非标准的非空听 $n_0$ 面听牌型。”
  \textit{Source: half-formal.qmd, Lemma \texttt{\#lem-bound}.}
  If there is a counterexample at some \(n\geq6\) divisible by \(3\),
  then there is an \(n_0\in\mathcal S\) at which a counterexample exists.
\end{theorem}
% SOURCE QUOTE PROOF: “假设反例集合非空。考虑所有出现反例的规模 $$\mathcal N=\left\{n\in\mathbb N:n\ge6,\ 3\mid n,\ \text{存在一个规模为 $n$ 的非标准全听牌型}\right\}.$$ 由自然数的良序性，$\mathcal N$ 有最小元素。记这个最小规模为 $n$，并固定一个规模为 $n$ 的非标准全听牌型 $c=(c_1,\ldots,c_n)$。如果 $n=6$，那么结论已经成立。以下设 $n>6$。由于 $3\mid n$，必有 $n\ge9$，所以 @lem-delete 可以用于 $c$。由 @lem-interior-positive，对每个内部坐标 $i\in\{2,\ldots,n-1\}$，都有 $1\le c_i\le3$。我们断言，内部序列 $c_2,c_3,\ldots,c_{n-1}$ 不含三个连续的 $1$。反设存在 $j$ 满足 $2\le j\le n-3$ 且 $c_j=c_{j+1}=c_{j+2}=1$。由 @lem-delete，删除这三个坐标得到一个非空听的 $(n-3)$ 面听牌型 $c'$。又因为 $n-3\ge6$、$3\mid(n-3)$，而 $n$ 是出现反例的最小规模，所以 $c'$ 不可能是非标准牌型，即 $c'=(3,1,\ldots,1,3)$。但 @lem-delete 的最后一个等价性随即推出 $c=(3,1,\ldots,1,3)$，与 $c$ 非标准矛盾。因此内部确实不存在三个连续的 $1$。下面估计内部长度。定义内部异常坐标集合 $E=\{i\in\{2,\ldots,n-1\}:c_i\ne1\}$，并令 $k=\#E$。由内部正性，对每个内部坐标都有 $c_i-1\ge0$；而当 $i\in E$ 时，$c_i\in\{2,3\}$，所以 $c_i-1\ge1$。因此 $$k\le\sum_{i=2}^{n-1}(c_i-1). \tag{10}$$ 利用总牌数条件计算右侧： $$\begin{aligned} \sum_{i=2}^{n-1}(c_i-1) &=\sum_{i=2}^{n-1}c_i-(n-2)\\ &=\left(\sum_{i=1}^{n}c_i-c_1-c_n\right)-(n-2)\\ &=(n+4-c_1-c_n)-(n-2)\\ &=6-c_1-c_n\\ &\le6. \end{aligned}$$ 最后一步使用 $c_1,c_n\ge0$。结合式 (10)，得到 $k\le6$。将内部序列按 $k$ 个异常坐标切开。若这些异常坐标依次为 $i_1<i_2<\cdots<i_k$，则内部的 $1$ 被分成至多 $k+1$ 段： 1. $i_1$ 之前的一段； 2. 每两个相邻异常坐标之间的一段，共 $k-1$ 段； 3. $i_k$ 之后的一段。当 $k=0$ 时，整个内部序列本身视为唯一的一段。因为已经证明不存在三个连续的 $1$，每一段的长度都至多为 $2$。所以内部坐标总数满足 $$n-2\le\underbrace{k}_{\text{异常坐标}}+\underbrace{2(k+1)}_{\text{至多 }k+1\text{ 段，每段至多两个 }1}=3k+2.$$ 由于 $k\le6$，可得 $n-2\le3\cdot6+2=20$，从而 $n\le22$。最后结合 $n\ge6$、$3\mid n$ 与 $n\le22$，可知 $n\in\{6,9,12,15,18,21\}$。我们选取的最小反例本身已经落在这个集合中，因此结论成立。”
\begin{proof}

\leanok
  Choose the least size \(n\geq6\), divisible by \(3\), at which a
  counterexample \(c\) exists.  If \(n=6\), there is nothing to prove.
  Otherwise \(n\geq9\).

  Every internal coordinate lies between \(1\) and \(3\) by
  Theorem~\ref{thm:interior_positive}.  There cannot be three consecutive
  internal coordinates equal to \(1\): deleting such a block would give a
  live all-rank wait of size \(n-3\).  Minimality would force the shortened
  hand to be standard, and
  Theorem~\ref{thm:delete_three_standard} would then force \(c\) itself to
  be standard, a contradiction.

  Let \(E\) be the set of internal coordinates whose multiplicity is not
  \(1\), and put \(k=|E|\).  Positivity and the upper bound \(3\) imply
  that every member of \(E\) contributes at least \(1\) to
  \(\sum_{i=2}^{n-1}(c_i-1)\).  Hence
  \[
    k\leq\sum_{i=2}^{n-1}(c_i-1)
      =(n+4-c_1-c_n)-(n-2)
      =6-c_1-c_n\leq6.
  \]
  The \(k\) exceptional positions divide the remaining \(1\)'s into at
  most \(k+1\) blocks.  Since no block has length \(3\), each has length
  at most \(2\), and therefore
  \[
    n-2\leq k+2(k+1)=3k+2\leq20.
  \]
  Thus \(n\leq22\).  Combining \(6\leq n\), \(3\mid n\), and
  \(n\leq22\) yields \(n\in\{6,9,12,15,18,21\}\), as required.
\end{proof}

\begin{corollary}\leanok
[Small-size exclusion implies global exclusion]
  \label{cor:no_counterexample}
  \lean{Ninegates.no_counterexample_of_small}
  \uses{def:counterexample_at,def:small_sizes,thm:counterexample_size_bound}
  % SOURCE: half-formal.qmd, equivalent form following #lem-bound (read from references/half-formal.qmd)
  % SOURCE QUOTE: “等价地说：若在所有 $n\in\{6,9,12,15,18,21\}$ 中都不存在反例，则在任何满足 $n\ge6$ 且 $3\mid n$ 的规模中都不存在反例。”
  \textit{Source: half-formal.qmd, equivalent form following Lemma
  \texttt{\#lem-bound}.}
  If no size in \(\mathcal S\) admits a counterexample, then no
  \(n\geq6\) divisible by \(3\) admits one.
\end{corollary}
% SOURCE QUOTE PROOF: “等价地说：若在所有 $n\in\{6,9,12,15,18,21\}$ 中都不存在反例，则在任何满足 $n\ge6$ 且 $3\mid n$ 的规模中都不存在反例。”
\begin{proof}

\leanok
  If a counterexample existed at an admissible size, apply
  Theorem~\ref{thm:counterexample_size_bound} to obtain one at a size in
  \(\mathcal S\), contradicting the hypothesis.
\end{proof}
